\subsection{Rings and Ideals}

Rings add multiplication to an abelian group. Ideals are the subsets by which quotient rings can be formed, just as subspaces are the subsets by which quotient vector spaces can be formed.

\begin{definition} \label{def:ring-commutative-identity-field} \hlt{Rings, Identities, and Fields}\\
Let $R$ be a set with operations $+$ and $\cdot$.
\begin{enumerate}[label=(\roman*)]
\setlength{\itemsep}{0pt}
\item $R$ is a \hlt{ring} if $(R,+)$ is an abelian group, multiplication is associative,\\
and $a(b+c)=ab+ac$, $(a+b)c=ac+bc$ for all $a,b,c\in R$.
\item A ring $R$ is \hlt{commutative} if $ab=ba$ for all $a,b\in R$.
\item A ring $R$ has \hlt{identity} if there is $1_R\in R$ such that $1_Ra=a1_R=a$ for all $a\in R$.
\item A \hlt{field} is a commutative ring with identity $1_R\neq0$ such that every $a\in R\setminus\{0\}$ has $a^{-1}\in R$.
\end{enumerate}
\end{definition}

\begin{example} \label{ex:basic-rings} \hlt{Basic Rings}\\
\begin{enumerate}[label=(\roman*)]
\setlength{\itemsep}{0pt}
\item $(\Z,+,\cdot)$ and $(\mathbb{F}[x],+,\cdot)$ are commutative rings with identity $1$.
\item $\operatorname{Mat}(n,n,\mathbb{F})$ is a unital ring under matrix addition and multiplication, with identity $I_n$;\\
it is noncommutative for $n\geq2$.
\item $2\Z$ is a commutative ring under inherited addition and multiplication, but has no identity.
\item $\Z/n\Z$ is a commutative ring under addition and multiplication modulo $n$, with identity $1+n\Z$.
\end{enumerate}
\end{example}

\begin{definition} \label{def:ring-homomorphism-and-subring} \hlt{Ring Homomorphisms and Subrings}\\
Let $R,S$ be rings with identity.
\begin{enumerate}[label=(\roman*)]
\setlength{\itemsep}{0pt}
\item A \hlt{ring homomorphism} $\varphi:R\to S$ preserves $+$, $\cdot$, and $1$, where\\
$\varphi(a+b)=\varphi(a)+\varphi(b)$, $\varphi(ab)=\varphi(a)\varphi(b)$, $\varphi(1_R)=1_S$ for all $a,b\in R$.
\item A \hlt{subring} of $R$ is a subset $T\subseteq R$ that is a ring under the inherited operations and has identity $1_R$.
\end{enumerate}
\end{definition}

\begin{proposition} \label{prop:subring-test} \hlt{Subring Test}\\
Let $R$ be a ring with identity and let $S\subseteq R$. Then $S\leq R \iff 1_R\in S$, $a-b\in S$, $ab\in S$ $\forall a,b\in S$.
\begin{proof}
\begin{enumerate}[label=\roman*.]
\setlength{\itemsep}{0pt}
\item[$(\Rightarrow)$] If $S\leq R$, then $1_R\in S$ and $S$ is closed under subtraction and multiplication.
\item[$(\Leftarrow)$] Assume $1_R\in S$, $a-b\in S$, and $ab\in S$ for all $a,b\in S$. Then $S$ is nonempty, $(S,+)\leq(R,+)$ by Proposition \ref{prop:subgroup-test}, and multiplication restricts to $S$. All ring laws and the identity $1_R$ are inherited from $R$, so $S\leq R$.
\end{enumerate}
\hlt{Structural Commentary}: A subring is a subset stable under the operations and containing the ambient unit.\\
\hlt{Key Insight}: Subtraction handles the additive group part in one stroke.
\end{proof}
\end{proposition}

\begin{definition} \label{def:ideal-and-generated-ideal} \label{ex:principal-ideal-polynomial-multiples} \hlt{Ideals and Generated Ideals}\\
Let $R$ be a ring.
\begin{enumerate}[label=(\roman*)]
\setlength{\itemsep}{0pt}
\item A nonempty subset $\mathcal{I}\subseteq R$ is a \hlt{two-sided ideal} if $a-b\in\mathcal{I}$ $\forall a,b\in\mathcal{I}$, and $ra,ar\in\mathcal{I}$ $\forall a\in\mathcal{I}$, $r\in R$.
\item If $R$ is commutative with identity, $S\subseteq R$, \hlt{ideal generated by $S$} is $\langle S\rangle\defeq\left\{\sum\limits_{j=1}^{n}r_js_j:n\geq1,\ r_j\in R,\ s_j\in S\right\}$.
\item \hlt{Finite generated ideal notation}: if $S=\{s_1,\ldots,s_n\}$, write $\langle s_1,\ldots,s_n\rangle\defeq\langle S\rangle$.
\item In $\mathbb{F}[x]$, $\langle p(x)\rangle=\{q(x)p(x):q(x)\in\mathbb{F}[x]\}$ is the \hlt{principal} ideal generated by $p(x)$.
\end{enumerate}
\end{definition}

\begin{proposition} \label{prop:generated-ideal-smallest-containing-set} \hlt{Generated Ideals Are Smallest}\\
Let $R$ be a commutative ring with identity and let $S\subseteq R$.\\
Then $\langle S\rangle$ is an ideal, $S\subseteq\langle S\rangle$, and for every ideal $\mathcal{J}\subseteq R$, $S\subseteq\mathcal{J}\implies \langle S\rangle\subseteq\mathcal{J}$.
\begin{proof}
Let $x=\sum_{j=1}^{m}r_js_j$ and $y=\sum_{k=1}^{n}t_ku_k$ lie in $\langle S\rangle$.\\
Then $x-y=\sum\limits_{j=1}^{m}r_js_j+\sum\limits_{k=1}^{n}(-t_k)u_k\in\langle S\rangle$; $ax=\sum\limits_{j=1}^{m}(ar_j)s_j\in\langle S\rangle$ for all $a\in R$, so $\langle S\rangle$ is an ideal.\\
Since $s=1_Rs\in\langle S\rangle$, $S\subseteq\langle S\rangle$. If $\mathcal{J}$ is an ideal with $S\subseteq\mathcal{J}$, then $r_js_j\in\mathcal{J}$ and finite sums remain in $\mathcal{J}$, QED.\\
\hlt{Structural Commentary}: Generated ideals close a set under exactly the operations ideals must allow.\\
\hlt{Key Insight}: To force $S$ into an ideal, one must include all ring-linear combinations of elements of $S$.
\end{proof}
\end{proposition}

\begin{proposition} \label{prop:intersections-and-chain-unions-of-ideals} \hlt{Intersections and Chain Unions of Ideals}\\
Let $R$ be a ring.
\begin{enumerate}[label=(\roman*)]
\setlength{\itemsep}{0pt}
\item If $\{\mathcal{I}_{\alpha}\}_{\alpha\in A}$ is a nonempty family of ideals of $R$, then $\bigcap_{\alpha\in A}\mathcal{I}_{\alpha}$ is an ideal of $R$.
\item If $\{\mathcal{I}_{\alpha}\}_{\alpha\in A}$ is a chain of ideals of $R$ under inclusion, then $\bigcup_{\alpha\in A}\mathcal{I}_{\alpha}$ is an ideal of $R$.
\end{enumerate}
\begin{proof}
\begin{enumerate}[label=(\roman*)]
\setlength{\itemsep}{0pt}
\item Let $\mathcal{I}=\bigcap_{\alpha\in A}\mathcal{I}_{\alpha}$. If $a,b\in\mathcal{I}$ and $r\in R$, then $a-b,ra,ar\in\mathcal{I}_{\alpha}$ for every $\alpha\in A$, hence $a-b,ra,ar\in\mathcal{I}$.
\item Let $\mathcal{U}=\bigcup_{\alpha\in A}\mathcal{I}_{\alpha}$. If $a,b\in\mathcal{U}$, then $a\in\mathcal{I}_{\alpha}$ and $b\in\mathcal{I}_{\beta}$ for some $\alpha,\beta$. Since the family is a chain, one of $\mathcal{I}_{\alpha},\mathcal{I}_{\beta}$ contains both $a,b$; hence $a-b\in\mathcal{U}$. If $a\in\mathcal{U}$ and $r\in R$, then $a\in\mathcal{I}_{\alpha}$ for some $\alpha$, so $ra,ar\in\mathcal{I}_{\alpha}\subseteq\mathcal{U}$.
\end{enumerate}
\hlt{Structural Commentary}: Chains avoid the usual failure of unions to preserve closure.\\
\hlt{Key Insight}: Comparability places any finite test inside one ideal.
\end{proof}
\end{proposition}

\begin{definition} \label{def:quotient-ring} \hlt{Quotient Rings}\\
Let $R$ be a ring and let $\mathcal{I}\trianglelefteq R$ be a two-sided ideal.
\begin{enumerate}[label=(\roman*)]
\setlength{\itemsep}{0pt}
\item The underlying set is the additive quotient group $R/\mathcal{I}=\{a+\mathcal{I}:a\in R\}$.
\item Addition and multiplication are $(a+\mathcal{I})+(b+\mathcal{I})=(a+b)+\mathcal{I}$ and $(a+\mathcal{I})(b+\mathcal{I})=ab+\mathcal{I}$.
\item With these operations, $R/\mathcal{I}$ is the \hlt{quotient ring} of $R$ by $\mathcal{I}$.
\end{enumerate}
\end{definition}

\begin{proposition} \label{prop:quotient-ring-well-defined-by-ideal} \hlt{Ideals Make Quotient Multiplication Well Defined}\\
Let $\mathcal{I}\leq (R,+)$. The rule $(a+\mathcal{I})(b+\mathcal{I})\defeq ab+\mathcal{I}$ is well defined on $R/\mathcal{I}$ iff $\mathcal{I}\trianglelefteq R$.
\begin{proof}
\begin{enumerate}[label=\roman*.]
\setlength{\itemsep}{0pt}
\item[$(\Rightarrow)$] If the rule is well defined and $i\in\mathcal{I}$, then $i+\mathcal{I}=0+\mathcal{I}$.\\
For $r\in R$, $ri+\mathcal{I}=r0+\mathcal{I}$, and $ir+\mathcal{I}=0r+\mathcal{I}$, so $ri,ir\in\mathcal{I}$. Thus $\mathcal{I}\trianglelefteq R$.
\item[$(\Leftarrow)$] If $\mathcal{I}\trianglelefteq R$ and $a'=a+i$, $b'=b+j$ with $i,j\in\mathcal{I}$, then $a'b'-ab=aj+ib+ij\in\mathcal{I}$.\\
Hence $a'b'+\mathcal{I}=ab+\mathcal{I}$, so the product is independent of representatives.
\end{enumerate}
\hlt{Structural Commentary}: Ideals are exactly the additive subgroups compatible with multiplication of cosets.\\
\hlt{Key Insight}: Coset multiplication is well defined exactly for ideals.
\end{proof}
\end{proposition}

\begin{definition} \label{def:maximal-ideal} \hlt{Maximal Ideals}\\
Let $R$ be a ring and let $\mathcal{M}\trianglelefteq R$.
\begin{enumerate}[label=(\roman*)]
\setlength{\itemsep}{0pt}
\item $\mathcal{M}$ is \hlt{proper} if $\mathcal{M}\neq R$.
\item $\mathcal{M}$ is \hlt{maximal} if $\mathcal{M}\neq R$ and, for every ideal $\mathcal{J}\trianglelefteq R$, $\mathcal{M}\subseteq\mathcal{J}\subseteq R \implies \mathcal{J}=\mathcal{M}\ \text{or}\ \mathcal{J}=R$.
\end{enumerate}
\end{definition}

\begin{theorem} \label{thm:quotient-by-maximal-ideal-is-field} \hlt{Maximal Ideals and Fields}\\
Let $R$ be a commutative ring with identity and let $\mathcal{I}\trianglelefteq R$. Then $R/\mathcal{I}$ is a field $\iff$ $\mathcal{I}$ is maximal
\begin{proof}
\begin{enumerate}[label=\roman*.]
\setlength{\itemsep}{0pt}
\item[$(\Rightarrow)$] If $R/\mathcal{I}$ is a field and $\mathcal{I}\subsetneq\mathcal{J}\trianglelefteq R$, then $\mathcal{J}/\mathcal{I}$ is a nonzero ideal of $R/\mathcal{I}$.\\
Hence $\mathcal{J}/\mathcal{I}=R/\mathcal{I}$, so $\mathcal{J}=R$. Thus $\mathcal{I}$ is maximal.
\item[$(\Leftarrow)$] If $\mathcal{I}$ is maximal and $0\neq a+\mathcal{I}\in R/\mathcal{I}$, then $a\notin\mathcal{I}$.\\
The ideal $\mathcal{I}+\langle a\rangle$ strictly contains $\mathcal{I}$, so maximality gives $\mathcal{I}+\langle a\rangle=R$. Thus $1=i+ra$ for some $i\in\mathcal{I}$, $r\in R$, and $(r+\mathcal{I})(a+\mathcal{I})=1+\mathcal{I}$. Hence every nonzero coset is invertible.
\end{enumerate}
\hlt{Structural Commentary}: Maximal ideals are exactly kernels whose quotients have no residual ideal structure.\\
\hlt{Key Insight}: Maximality makes every nonzero coset invertible.
\end{proof}
\end{theorem}

\begin{theorem} \label{thm:maximal-ideals-exist} \hlt{Existence of Maximal Ideals}\\
If $R$ is a commutative ring with identity and $R\neq0$, then $R$ has a maximal ideal.
\begin{proof}
Let $\mathcal{P}\defeq\{\mathcal{I}\trianglelefteq R:\mathcal{I}\neq R\}$, ordered by inclusion. Since $R\neq0$, $\{0\}\in\mathcal{P}$.\\
If $\mathcal{C}\subseteq\mathcal{P}$ is a chain, then $\mathcal{U}\defeq\bigcup_{\mathcal{I}\in\mathcal{C}}\mathcal{I}$ is an ideal by Proposition \ref{prop:intersections-and-chain-unions-of-ideals}.\\
Also $\mathcal{U}\neq R$: if $1_R\in\mathcal{U}$, then $1_R\in\mathcal{I}$ for some $\mathcal{I}\in\mathcal{C}$, forcing $\mathcal{I}=R$, contradiction. Thus $\mathcal{U}\in\mathcal{P}$ is an upper bound for $\mathcal{C}$. By Zorn's lemma, Principle \ref{princ:zorns-lemma}, $\mathcal{P}$ has a maximal element, i.e. $R$ has a maximal ideal.\\
\hlt{Structural Commentary}: Maximal ideals are produced by maximality, not by an explicit formula in general.\\
\hlt{Key Insight}: Chain unions keep properness because containing $1$ is a finite obstruction.
\end{proof}
\end{theorem}

\begin{definition} \label{def:zero-divisor-integral-domain} \hlt{Zero Divisors and Integral Domains}\\
Let $R$ be a ring.
\begin{enumerate}[label=(\roman*)]
\setlength{\itemsep}{0pt}
\item A \hlt{zero divisor} is a nonzero $a\in R$ such that $\exists\,0\neq b\in R$ with $ab=0$ or $ba=0$.
\item An \hlt{integral domain} is a commutative ring with identity $1_R\neq0$ and no zero divisors.
\end{enumerate}
\end{definition}

\begin{proposition} \label{prop:integral-domain-cancellation} \hlt{Cancellation in Integral Domains}\\
Let $R$ be a commutative ring with identity $1_R\neq0$.\\
Then $R$ is an integral domain $\iff$ $(\forall r,x,y\in R)(r\neq0\land rx=ry\Rightarrow x=y)$.
\begin{proof}
\begin{enumerate}[label=\roman*.]
\setlength{\itemsep}{0pt}
\item[$(\Rightarrow)$] If $rx=ry$ and $r\neq0$, then $r(x-y)=0$; no zero divisors gives $x-y=0$.
\item[$(\Leftarrow)$] If $ab=0$ and $a\neq0$, then $ab=a0$; cancellation gives $b=0$, so no zero divisors exist.
\end{enumerate}
\hlt{Structural Commentary}: Domains are exactly cancellative commutative rings.\\
\hlt{Key Insight}: Zero divisors obstruct cancellation.
\end{proof}
\end{proposition}

\begin{definition} \label{def:field-of-quotients} \hlt{Field of Quotients}\\
Let $R$ be an integral domain. On $R\times(R\setminus\{0\})$, set $(a,b)\sim(c,d)$ $\Longleftrightarrow$ $ad=bc$.\\
The \hlt{field of quotients} of $R$ is $(R\times(R\setminus\{0\}))/\sim$; write $[(a,b)]=a/b$. For $b,d\neq0$, $\frac{a}{b}+\frac{c}{d}=\frac{ad+bc}{bd}$; $\frac{a}{b}\frac{c}{d}=\frac{ac}{bd}$.
\end{definition}

\begin{proposition} \label{prop:field-of-quotients-embeds-domain} \hlt{Domains Embed in Fields}\\
Let $R$ be an integral domain with field of quotients $K$.\\
Then $\iota:R\hookrightarrow K$; $\iota(a)=a/1$  is an injective ring homomorphism.
\begin{proof}
By Proposition \ref{prop:integral-domain-cancellation}, the relation and operations in Definition \ref{def:field-of-quotients} are well defined.\\
For $a,b\in R$, $\iota(a+b)=\frac{a+b}{1}=\frac{a}{1}+\frac{b}{1}$; $\iota(ab)=\frac{ab}{1}=\frac{a}{1}\frac{b}{1}$; and $\iota(1_R)=\frac{1_R}{1}$..\\
Thus $\iota$ is a ring homomorphism. If $\iota(a)=\iota(b)$, then $a/1=b/1$, so $a=b$; hence $\iota$ is injective.\\
\hlt{Structural Commentary}: The field of quotients formally adds denominators to a domain.\\
\hlt{Key Insight}: Integral domains are exactly the setting where fraction arithmetic is consistent.
\end{proof}
\end{proposition}

\begin{definition} \label{def:characteristic-of-ring} \hlt{Characteristic}\\
Let $R$ be a ring with identity. Its \hlt{characteristic} is $\operatorname{char}(R)\defeq
\begin{cases}
\min\{n\in\mathbb{N}_{>0}:n1_R=0_R\}, & \text{if this set is nonempty},\\
0, & \text{otherwise}.
\end{cases}$.
\end{definition}

\begin{proposition} \label{prop:finite-domain-prime-characteristic} \hlt{Finite Domains Have Prime Characteristic}\\
Let $R$ be a finite ring with identity. Then $\operatorname{char}(R)\neq0$. If $R$ is an integral domain, then $\operatorname{char}(R)$ is prime.
\begin{proof}
Since $R$ is finite, the map $n\mapsto n1_R$ from $\mathbb{N}$ to $R$ is not injective.\\
Thus $i1_R=j1_R$ for some $i<j$, so $(j-i)1_R=0_R$ and $\operatorname{char}(R)>0$.\\
Assume $R$ is a domain and set $c=\operatorname{char}(R)$. Since $1_R\neq0_R$, $c>1$.\\
If $c=ab$ with $1<a,b<c$, then $(a1_R)(b1_R)=ab1_R=c1_R=0_R$.\\
By minimality of $c$, $a1_R\neq0_R$ and $b1_R\neq0_R$, contradicting that $R$ has no zero divisors. Hence $c$ is prime.\\
\hlt{Structural Commentary}: Characteristic records the additive order of the unit.\\
\hlt{Key Insight}: In a domain, a composite characteristic would factor zero nontrivially.
\end{proof}
\end{proposition}

\begin{definition} \label{def:algebra-over-field} \hlt{Algebras over a Field}\\
Let $\mathbb{F}$ be a field.\\
An \hlt{$\mathbb{F}$-algebra} is an $\mathbb{F}$-vector space $\mathcal{A}$ with a ring product satisfying $\lambda(ab)=(\lambda a)b=a(\lambda b)$ $\forall \lambda\in\mathbb{F}$ and $a,b\in\mathcal{A}$.
\end{definition}

\begin{example} \label{ex:matrix-and-polynomial-algebras} \hlt{Matrix and Polynomial Algebras}\\
$\operatorname{Mat}(n,n,\mathbb{F})$ and $\mathbb{F}[x]$ are algebras over $\mathbb{F}$. The first is generally noncommutative; the second is commutative.
\end{example}
